\chapter{Background and Related Work}
\label{ch:related}

Our work sits at the intersection of three research threads: policy distillation, memory architectures for sequential decision-making, and learning from imperfect demonstrations.
Each has developed largely independently; our contribution is to study how they interact.

\paragraph{Policy distillation.}
The idea of compressing a large model's knowledge into a smaller one originated with~\citet{hinton2015distilling}, who showed that a teacher's soft output distribution carries richer signal than hard labels.
\citet{rusu2016policy} brought this to RL, distilling DQN policies into compact students, and~\citet{parisotto2016actormimic} extended it to multi-task transfer.
\citet{czarnecki2019distilling} studied when and why distillation fails, finding that objective choice and capacity mismatch are key failure modes.
More recently, \citet{busbridge2025distillation} studied scaling properties of distillation in language models, finding teacher quality to be a strong predictor of student quality---a result that motivates our central question: does this finding extend to sequential decision-making under partial observability, where the student must additionally handle temporal memory?
\citet{belkhale2023data} showed that data quality dominates other design factors in imitation learning for manipulation, providing further evidence that the quality of demonstrations may matter more than architectural sophistication.
Our work extends this line of inquiry to POMDPs with a systematic architecture comparison and formal variance decomposition.

\paragraph{Memory architectures for RL.}
The question of which temporal architecture best handles partial observability has been studied extensively in the RL training setting.
The POPGym benchmark~\citep{morad2023popgym} compared 13 architectures across 15 POMDP tasks and found GRU to be the strongest general-purpose choice---a finding we corroborate in the distillation setting.
\citet{ni2022transformers} showed that simple recurrent baselines are surprisingly competitive with Transformers on many POMDPs, challenging the assumption that attention-based models dominate.
\citet{chen2021decision} reframed RL as sequence prediction with Decision Transformer, establishing that Transformers can be effective for sequential decision-making.
However, all these studies compare architectures trained \emph{via RL}, where the architecture's inductive bias shapes exploration and credit assignment.
Distillation eliminates both---the teacher has already explored, and targets are available at every timestep---raising the question of whether RL-derived architecture rankings transfer to the distillation setting.

\paragraph{Spectral and state space models.}
A parallel line of work has developed structured sequence models that replace learned recurrence with mathematically-principled temporal processing.
\citet{hazan2017learning} introduced spectral filtering for online prediction of linear dynamical systems.
\citet{agarwal2024spectral} formalized the Spectral Transform Unit (STU) with fixed eigenvectors of a Hankel matrix and convexity guarantees for the learnable parameters.
\citet{liu2025flashstu} proposed Flash STU with a tensordot approximation for parameter efficiency at large scale.
S4~\citep{gu2022s4} and Mamba~\citep{gu2024mamba} take a different approach, learning their state-transition dynamics via structured parameterizations.
We focus on STU rather than Mamba because we aim to test the \emph{fixed} spectral prior against learned dynamics (GRU/LSTM)---a cleaner comparison than learned-SSM vs learned-RNN, where both sides learn their temporal basis.
Our discovery that the Flash STU tensordot approximation fails catastrophically at compact model scales is, to our knowledge, unreported in prior work.

\paragraph{Learning from imperfect demonstrations.}
Our finding that students frequently exceed their imperfect teachers connects to a broader literature on learning from suboptimal demonstrations.
\citet{ross2011dagger} introduced DAgger to address the distribution shift problem in behavioral cloning---the fundamental limitation that offline-collected demonstrations may not cover the states the student encounters.
\citet{beliaev2022ileed} showed that explicitly modeling demonstrator expertise allows learning from mixed-quality demonstrations, and~\citet{emmons2022rvs} found that simple MLPs often match fancier architectures in offline supervised RL---a finding our Ant baseline results echo.
Our proposed mechanism (variance averaging + length-proportional loss weighting) adds a new hypothesis to this literature, though we note it requires causal testing to confirm.
